\documentclass[11pt]{article}
\usepackage[a4paper,margin=26mm]{geometry}
\usepackage[T1]{fontenc}
\usepackage{lmodern}
\usepackage{amsmath,amssymb,amsthm}
\usepackage{microtype}
\usepackage{booktabs}
\usepackage{array}
\usepackage{xurl}
\usepackage{hyperref}
\hypersetup{colorlinks=true,linkcolor=blue,urlcolor=blue,citecolor=blue,
 pdftitle={The minimum order of a counterexample to a conjecture on integral generalized sun graphs},
 pdfauthor={Horace Dong}}
\newtheorem{lemma}{Lemma}[section]
\newtheorem{corollary}[lemma]{Corollary}
\newtheorem*{thmone}{Theorem 1}
\newtheorem*{thmtwo}{Theorem 2}
\newtheorem*{propthree}{Proposition 3}
\newtheorem*{conjtwo}{Conjecture 2}
\theoremstyle{remark}
\newtheorem{remark}[lemma]{Remark}
\DeclareMathOperator{\mult}{mult}
\DeclareMathOperator{\tr}{tr}
\DeclareMathOperator{\diag}{diag}
\newcommand{\Z}{\mathbb Z}
\newcommand{\R}{\mathbb R}
\DeclareUrlCommand\code{\urlstyle{tt}\def\UrlBreaks{\do\_\do\/\do\.}\def\UrlBigBreaks{}}
\setlength{\emergencystretch}{1.5em}
\title{The minimum order of a counterexample to a conjecture\\ on integral generalized sun graphs}
\author{Horace Dong\\ \small Independent researcher\\
\small\texttt{maxwelldhx@gmail.com}\\
\small ORCID: \href{https://orcid.org/0009-0002-7554-7548}{0009-0002-7554-7548}}
\date{September 2026}
\begin{document}
\maketitle
\begin{abstract}
A graph is integral if its adjacency eigenvalues are integers, and a
generalized sun graph is a cycle with pendant paths attached to its vertices.
Braga, Del-Vecchio and Rodrigues conjectured that the
cycle of an integral generalized sun graph that is not a cycle has length
divisible by $4$. Braga, Moraes and Santos disproved this with the graph
$C_{6,1}(0,6,6,12,6,6)$ on $42$ vertices, and remarked that their computer
search does not show that no smaller counterexample exists. We show that
none exists: up to isomorphism, $C_{6,1}(0,6,6,12,6,6)$ is the only
counterexample with at most $42$ vertices, also when pendant paths of both
admissible lengths share a vertex. Every integral generalized sun graph with
at most $41$ vertices that is not a cycle has a cycle of length $4$ and is one
of the three graphs found by Braga, Rodrigues and Trevisan in 2017. Up to $42$
vertices, spectral lemmas exclude every cycle length $b\not\equiv0\pmod 4$ except $3$, $6$ and
$10$, which are settled by exhaustive searches with separately written programs,
two of them by independent referee agents. We
also note that if an integral unicyclic graph has a cycle of
odd length $b$, then every prime $q$ with $(q-1)\mid(b-1)$ divides $2b$;
hence $3\mid b$, and $b=3$ or $b\ge15$. Nothing in this paper is formally
verified. The results were found by AI coding agents working under the
author's direction; their role is described at the end of the paper.
\end{abstract}

\section{Introduction}

A graph is \emph{integral} if all eigenvalues of its adjacency matrix are
integers, a notion that goes back to Harary and Schwenk~\cite{hs}. A search of
Braga, Rodrigues and Trevisan~\cite{brt} showed that, besides the cycles $C_3$,
$C_4$ and $C_6$, exactly three unicyclic graphs with at most $21$ vertices are
integral. Braga, Del-Vecchio and Rodrigues~\cite{bdr} constructed infinite
families of integral unicyclic graphs by attaching pendant paths to the
vertices of a cycle; such graphs are called \emph{generalized sun graphs}.
Following~\cite{bms}, $C_b(n_1P_{t_1},\dots,n_bP_{t_b})$ is obtained from the
cycle $v_1\cdots v_bv_1$ by attaching $n_k$ copies of the path $P_{t_k}$ on
$t_k$ vertices to $v_k$ (joining $v_k$ to an end of each copy), and
$C_{b,t}(n_1,\dots,n_b)$ is the case $t_1=\dots=t_b=t$. We state the second
conjecture of~\cite{bdr} in our words.

\begin{conjtwo}[{\cite{bdr}; see~\cite[Conjecture~2]{bms}}]
Let $G$ be an integral generalized sun graph that is not a cycle. Then the
cycle of $G$ has length divisible by $4$.
\end{conjtwo}

Braga, Moraes and Santos~\cite{bms} proved the first conjecture of~\cite{bdr}
in a strong form: a graph with a pendant path of at least three edges has an
eigenvalue in $(1,2\cos(\pi/9)]$, so it is not integral~\cite[Theorem~2.1]{bms}.
They disproved Conjecture~2: $C_{6,1}(0,6,6,12,6,6)$ has $42$ vertices and
spectrum $\{0^{[32]},\pm2^{[2]},\pm3^{[2]},\pm4\}$~\cite[Proposition~3.2]{bms},
and it is the smallest member of an infinite
family~\cite[Corollary~3.6 and the remark after it]{bms}.
Their search over boxes of parameter vectors, for graphs with pendant vertices
only, found no other integral graph whose cycle length is not divisible by
$4$. As they point out~\cite[Section~4]{bms}, it ``does not prove that the
graph on 42 vertices has minimum order'' among all counterexamples, since
vectors outside the boxes may give smaller graphs, and graphs with copies of
$P_2$ were not searched. This note answers that question.

A \emph{counterexample} is a generalized sun graph that is integral, is not a
cycle, and has cycle length $b\not\equiv0\pmod4$. We allow any finite number of
pendant paths at each cycle vertex, possibly of different lengths at the same
vertex. By~\cite[Theorem~2.1]{bms}, an integral generalized sun graph
then carries only copies of $P_1$ and $P_2$. The graphs
$C_b(n_1P_{t_1},\dots,n_bP_{t_b})$ form a subclass containing every graph
named below, so our results hold in either reading.

\begin{thmone}
Up to isomorphism, $C_{6,1}(0,6,6,12,6,6)$ is the only counterexample to
Conjecture~2 with at most $42$ vertices. In particular, the minimum order of a
counterexample is $42$.
\end{thmone}

\begin{thmtwo}
Every integral generalized sun graph with at most $41$ vertices that is not a
cycle has a cycle of length $4$, and is isomorphic to $C_{4,1}(6,0,3,0)$, with
$13$ vertices, or to $C_4(5P_2,P_1,2P_2,P_1)$ or $C_{4,2}(4,4,0,0)$, with $20$
vertices each.
\end{thmtwo}

These three graphs are the three integral unicyclic graphs found by Braga,
Rodrigues and Trevisan~\cite[Section~5, p.~296, and Table~1, p.~297]{brt}, with
spectra $\{0^{[9]},\pm2,\pm3\}$ and, for the two cospectral graphs,
$\{0^{[2]},\pm1^{[7]},\pm2,\pm3\}$. Theorem~2 adds that there is no other
integral generalized sun graph up to $41$ vertices. The one-page abstract of
Veras~\cite{ver} reports a computer search for integral unicyclic graphs. It
names $C_{4,1}(0,16,32,16)$, on $68$ vertices, and the family
$C_{4,1}(0,a,2a,a)$, whose integral members have at least $68$
vertices~\cite{bms}. It does not state the range of the search and names no
graph with at most $41$ vertices, so nothing it reports overlaps Theorem~2,
although the search may have covered part of that range. Write $\mult(\lambda)$ for
the multiplicity of an eigenvalue $\lambda$. The proofs use the following
observation on integral graphs.

\begin{propthree}
Let $G$ be an integral graph with $e$ edges that contains an odd cycle, let $g$
be the length of its shortest odd cycles and $c_g$ their number. Then every
prime $q$ with $(q-1)\mid(g-1)$ divides $2g\,c_g$; in particular $3\mid g\,c_g$.
Moreover, at least $(g-1)/2$ integers $v\ge2$ satisfy $\mult(v)\ne\mult(-v)$,
so $2e\ge\sum_{v=2}^{(g+1)/2}v^2$.
\end{propthree}

For an integral unicyclic graph with a cycle of odd length $b$ this gives
$3\mid b$, and $b=3$ or $b\ge15$, with at least $102$ vertices in the latter
case (Corollary~\ref{cor:odd}).

\paragraph{Method.}
Up to $42$ vertices, hand proofs exclude odd $b\ge5$ (Proposition~3),
$b\in\{14,18\}$ (Corollary~\ref{cor:1418}) and even $b\ge22$
(Lemma~\ref{lem:cycle}), and a bound due to the first referee agent (see the
disclosure; Lemma~\ref{lem:rank}) excludes every $b\ge14$ again; exhaustive
searches settle $b=3,6,10$ (Section~\ref{sec:comp}). Every case is confirmed at
least twice, by hand proofs or by separately written programs
(Table~\ref{tab:cases}); two of the programs, C and D, were written from
scratch by referee agents that had not produced the results. Nothing here is
formally verified, in Lean or otherwise; the results rest on the written proofs
and the computations.

\paragraph{Prior work.}
The question was raised in version~1 of~\cite{bms}, submitted to the arXiv on
23 September 2026. We found no earlier answer to it and no discussion of it.
Our searches on 26 September 2026, 18:58--19:43 UTC, covered the arXiv listing
of~\cite{bms} (version~1 only); Crossref; MathDB
(\url{https://mathdb.com}); GitHub issues, pull requests, commits, code and
repositories, including the \texttt{trureturing} repository~\cite{tru} of
AI-assisted work on open problems (as a positive control, the same search for
the versioned identifier of another recent arXiv paper found the expected
issue); MathOverflow and Mathematics Stack Exchange; and web search. The
OpenAlex search and citation queries failed with HTTP error 429 (rate limit), and several GitHub
searches were rate-limited and repeated later. From 19:59 to 20:32 UTC the
first referee agent repeated part of these searches and added arXiv API queries and
full-text searches (Remark~\ref{rem:novelty}). We did not search Google
Scholar, zbMATH or MathSciNet. The earlier computer searches for integral
unicyclic graphs of~\cite{brt} and~\cite{ver} are discussed after Theorem~2;
we read the abstract~\cite{ver}, which~\cite{bms} cite, on 26 September 2026 at
22:48 UTC, and at 22:56 UTC the arXiv still listed only version~1
of~\cite{bms}. A final check on 26 September 2026 between 23:15 and 23:17 UTC
(the arXiv, GitHub, MathDB and MathOverflow) found nothing new. This reports
the coverage of our searches, not a guarantee of priority.

\section{Notation and spectral lemmas}\label{sec:lemmas}

Throughout, $G$ is a generalized sun graph with cycle $v_0v_1\cdots v_{b-1}v_0$
($b\ge3$, indices modulo $b$) carrying only copies of $P_1$ and $P_2$: at $v_k$
hang $p_k$ pendant vertices and $q_k$ pendant paths $v_kww'$. Put
$P=\sum p_k$, $Q=\sum q_k$, $K_p=\{k:p_k>0\}$, $K_q=\{k:q_k>0\}$,
$Z=\Z_b\setminus K_q$ and $A=A(G)$. Then $G$ has $n=b+P+2Q$ vertices and $n$
edges. Let $\Theta(G)=\sum_vd_v(d_v-1)$ over all vertices, with degrees $d_v$;
then $\Theta(G)=\sum_k\bigl((2+p_k+q_k)(1+p_k+q_k)+2q_k\bigr)$. An isomorphism
maps the unique cycle onto the unique cycle, and $(p_k,q_k)$ is determined by
the tree attached at $v_k$, so isomorphism classes correspond to dihedral
orbits of the sequences $((p_k,q_k))_{k\in\Z_b}$.

For a real symmetric matrix $M$, $n_+(M)$, $n_-(M)$ and $n_0(M)$ count its
positive, negative and zero eigenvalues. For a graph $H$, $N_{>t}(H)$,
$N_{\ge t}(H)$, $N_{<t}(H)$ and $N_{\le t}(H)$ count the eigenvalues of $A(H)$
in the indicated ranges, with multiplicity. If $G$ is integral, $B^+$ is the
multiset of its eigenvalues $\lambda\ge2$ and $S_2=\sum_{\lambda\in
B^+}\lambda^2$. For nonempty $K=\{k_1<\dots<k_t\}\subseteq\{0,\dots,b-1\}$ put
$k_{t+1}=k_1+b$; the \emph{gaps} of $K$ are the pairs $(k_i,k_{i+1})$, of
\emph{length} $k_{i+1}-k_i$. If $|K|=1$, the only gap has length $b$.

\begin{lemma}\label{lem:schur}
For $\lambda\in\R\setminus\{-1,0,1\}$ let
$f_k(\lambda)=p_k/\lambda+q_k\lambda/(\lambda^2-1)$ and
$S(\lambda)=A(C_b)-\lambda I+\diag(f_k(\lambda))_k$.
\begin{itemize}
\item[(a)] If $\lambda>1$, then $N_{>\lambda}(G)=n_+(S(\lambda))$ and
 $N_{\ge\lambda}(G)=n_+(S(\lambda))+n_0(S(\lambda))$.
\item[(b)] If $\lambda<-1$, then $N_{<\lambda}(G)=n_-(S(\lambda))$ and
 $N_{\le\lambda}(G)=n_-(S(\lambda))+n_0(S(\lambda))$.
\item[(c)] $N_{>1}(G)\ge|K_q|+n_+(R)$ and $N_{<-1}(G)\ge|K_q|+n_-(R')$, where
 $R$ and $R'$ are the principal submatrices of $A(C_b)-I+\diag(p_k)$ and
 $A(C_b)+I-\diag(p_k)$ on $Z$.
\end{itemize}
\end{lemma}

\begin{proof}
With the cycle first,
$A-\lambda I=\bigl(\begin{smallmatrix}A(C_b)-\lambda I&W\\W^{\mathsf
T}&H\end{smallmatrix}\bigr)$, where $H=H(\lambda)$ is the direct sum of a block
$(-\lambda)$ for each pendant vertex and a block
$\bigl(\begin{smallmatrix}-\lambda&1\\1&-\lambda\end{smallmatrix}\bigr)$ for
each pendant path (rows $w,w'$), and $W$ has one entry $1$ in the column of
each pendant vertex and each $w$, in the row of its neighbour on the cycle. $H$ has
eigenvalues $-\lambda$, $-\lambda\pm1$, so it is negative (positive) definite for
$\lambda>1$ ($\lambda<-1$). The $2\times2$ block has inverse with $(1,1)$ entry
$-\lambda/(\lambda^2-1)$, so $WH^{-1}W^{\mathsf T}=-\diag(f_k(\lambda))$ and
$A-\lambda I$ is congruent to $S(\lambda)\oplus H$; now use Sylvester's law of
inertia.

(c) As $\lambda\downarrow1$, the entries $f_k(\lambda)$, $k\in K_q$, tend to
$+\infty$ while the other entries of $S(\lambda)$ stay bounded; so the block
$S_{KK}$, $K=K_q$, is positive definite for $\lambda$ near $1$ and its inverse
tends to $0$. By Haynsworth's inertia additivity~\cite{hay},
$n_+(S(\lambda))=|K|+n_+(S_{ZZ}-S_{ZK}S_{KK}^{-1}S_{KZ})$, and the Schur
complement tends to $R$, so its positive inertia is eventually at least
$n_+(R)$. Finally $N_{>1}(G)=N_{>\lambda}(G)$ for $\lambda>1$ close to $1$. The
case $\lambda\uparrow-1$ is the same, with $f_k(\lambda)\to-\infty$ on $K_q$.
\end{proof}

\begin{lemma}\label{lem:mult}
\begin{itemize}
\item[(a)] Let $m\in\R\setminus\{-1,0,1\}$, $c_k=m-f_k(m)$,
 $T_k=\bigl(\begin{smallmatrix}c_k&-1\\1&0\end{smallmatrix}\bigr)$ and
 $M=T_{b-1}\cdots T_0$. Then $\mult(m)=\dim\ker(M-I)\le2$; it equals $2$ if
 and only if $M=I$, and $1$ if and only if $M\ne I$ and $\tr M=2$.
\item[(b)] $\mult(0)=P-|K_p|+\delta_0$, where, if $K_p=\emptyset$, $\delta_0=2$
 when $4\mid b$ and $\delta_0=0$ otherwise, and, if $K_p\ne\emptyset$,
 $\delta_0$ is the number of gaps of $K_p$ of even length.
\item[(c)] For $\varepsilon=\pm1$, $\mult(\varepsilon)=Q-|K_q|+\delta_\varepsilon$,
 where $\delta_\varepsilon$ is the dimension of the space of $y\in\R^{\Z_b}$
 with $y_k=0$ for $k\in K_q$ and $y_{k-1}+y_{k+1}=\varepsilon(1-p_k)y_k$ for
 $k\notin K_q$. Each gap of $K_q$ contributes $0$ or $1$ to
 $\delta_\varepsilon$, and $\delta_\varepsilon\le2$ if $K_q=\emptyset$.
\item[(d)] $\mult(0)\le\max(P,2)$ and $\mult(\pm1)\le\max(Q,2)$.
\end{itemize}
\end{lemma}

\begin{proof}
(a) If $Ax=mx$, then $x_a=x_{v_k}/m$ at a pendant vertex $a$ of $v_k$, and
$x_{w'}=x_w/m$, $x_w=mx_{v_k}/(m^2-1)$ on a pendant path $v_kww'$. So $x$ is
determined by its restriction $y$ to the cycle, which satisfies
$y_{k-1}+y_{k+1}=c_ky_k$, and every solution $y$ extends to an eigenvector.
Writing $(y_{k+1},y_k)^{\mathsf T}=T_k(y_k,y_{k-1})^{\mathsf T}$, the solutions
correspond to the vectors fixed by $M$; as $\det M=1$, $M-I$ is singular if and
only if $\tr M=2$.

(b) If $Ax=0$, a pendant vertex at $v_k$ forces $x_{v_k}=0$, and a pendant
path gives $x_w=0$, $x_{w'}=-x_{v_k}$. For $k\in K_p$ the equation at $v_k$
only fixes the sum of the $p_k$ pendant values ($p_k-1$ free values), and for
$k\notin K_p$ it reads $y_{k-1}+y_{k+1}=0$. So $\delta_0$ is the dimension of
the space of $y$ with $y_k=0$ on $K_p$ and $y_{k-1}+y_{k+1}=0$ off $K_p$. If
$K_p=\emptyset$, then $y_{k+2}=-y_k$ for all $k$, with a $2$-dimensional space
of $b$-periodic solutions if $4\mid b$ and none otherwise. If
$K_p\ne\emptyset$, the equations split along the gaps. On a gap of length $d$
put $z_j=y_{k_i+j}$; then $z_0=z_d=0$ and $z_{j-1}+z_{j+1}=0$ for $0<j<d$,
which forces $z_j=0$ for even $j$ and $z_{j+2}=-z_j$ for odd $j$. For even $d$
this leaves $z_1$ free; for odd $d$ it gives $z_1=\pm z_d=0$.

(c) If $Ax=\varepsilon x$, a pendant path gives $x_{w'}=\varepsilon x_w$ and
$\varepsilon x_w=x_{v_k}+x_{w'}$, so $x_{v_k}=0$ for $k\in K_q$, and the equation
at $v_k$ only fixes the sum of the $q_k$ values $x_w$ ($q_k-1$ free values). A pendant vertex gives
$x_a=\varepsilon x_{v_k}$, so for $k\notin K_q$ the equation at $v_k$ reads
$y_{k-1}+y_{k+1}+\varepsilon p_ky_k=\varepsilon y_k$. On a gap of $K_q$ the
recurrence determines all values from the first one after the left end, so
the gap contributes at most $1$; if $K_q=\emptyset$ the periodic solutions of a
recurrence of order two form a space of dimension at most $2$.

(d) If $K_p\ne\emptyset$, then $\delta_0\le|K_p|$, the number of gaps, so
$\mult(0)\le P$; otherwise $\mult(0)=\delta_0\le2$. Likewise for $\pm1$.
\end{proof}

\begin{lemma}\label{lem:trace}
Let $H$ be a graph with $e$ edges, degrees $d_v$, $c_4$ cycles of length $4$,
and adjacency matrix $A$.
\begin{itemize}
\item[(a)] $\tr A^2=2e$ and $\tr A^4=2e+2\sum_vd_v(d_v-1)+8c_4$.
\item[(b)] If the shortest odd cycles of $H$ have length $g$ and there are
 $c_g$ of them, then $\tr A^L=0$ for odd $L<g$ and $\tr A^g=2g\,c_g$.
\end{itemize}
\end{lemma}

\begin{proof}
We count closed walks. (a) A closed walk of length $4$ traverses one edge four
times ($2$ walks per edge), or two edges with a common vertex twice each ($4$
walks for each of the $\sum_v\binom{d_v}2$ such pairs), or a $4$-cycle ($8$ walks
per cycle). (b) A closed walk of odd length $L$ that repeats a vertex splits
there into two shorter closed walks, one of odd length. By induction it
contains an odd cycle of length at most $L$, with equality only if it
traverses a cycle once. So no closed walk of odd length $L<g$ exists, and each
$g$-cycle gives $2g$ closed walks of length $g$.
\end{proof}

For our graphs $e=n$, $c_4=[b=4]$, $\sum_vd_v(d_v-1)=\Theta(G)$, and $g=b$,
$c_g=1$ if $b$ is odd. If $G$ is integral, the eigenvalues in $\{0,\pm1\}$
satisfy $\lambda^4=\lambda^2$ and $\lambda^L=\lambda$ for odd $L$, and
$\tr A=0$, so
\begin{gather}
 \textstyle\sum_\lambda\lambda^2=2n,\qquad
 \sum_{|\lambda|\ge2}\lambda^2(\lambda^2-1)=2\,\Theta(G)+8\,[b=4],\label{eq:m4}\\
 \textstyle\sum_{|\lambda|\ge2}(\lambda^L-\lambda)=0\ \ (L\text{ odd},\ 3\le L<b),
 \qquad\sum_{|\lambda|\ge2}(\lambda^b-\lambda)=2b\ \ \text{if $b$ is odd}.\label{eq:odd}
\end{gather}

\begin{lemma}\label{lem:adm}
Let $G$ be integral and not a cycle.
\begin{itemize}
\item[(a)] The largest eigenvalue $\rho$ is simple and $\rho\ge3$; every
 eigenvalue satisfies $|\lambda|\le\rho$ and $|\lambda|\le\sqrt{2n}$, and
 $\lambda>-\rho$ if $b$ is odd.
\item[(b)] $\mult(\lambda)\le2$ whenever $|\lambda|\ge2$.
\item[(c)] If $b$ is even, the spectrum is symmetric, $S_2=n-\mult(1)$ and
 $P+Q\le S_2-b+2$. If $b$ is odd, $P+Q\le\frac12\sum_{|\lambda|\ge2}\lambda^2-b+2$.
\item[(d)] If $b$ is even, then $|B^+|\le4$ if $n\le41$, and $|B^+|\le5$ if
 $n\le42$, with equality only if $B^+=\{4,3,3,2,2\}$.
\end{itemize}
\end{lemma}

\begin{proof}
(a) By Perron--Frobenius theory~\cite{bh}, $\rho$ is simple, $|\lambda|\le\rho$,
and $-\rho$ is an eigenvalue only if $G$ is bipartite. As $G$ properly contains
$C_b$, $\rho>2$. Also $\lambda^2\le\tr A^2=2n$. (b) is
Lemma~\ref{lem:mult}(a). (c) For even $b$, $G$ is bipartite and~\eqref{eq:m4}
gives $2n=2S_2+2\mult(1)$; by Lemma~\ref{lem:mult}(d),
$S_2\ge b+P+2Q-\max(Q,2)\ge b+P+Q-2$. For odd $b$,
$\sum_{|\lambda|\ge2}\lambda^2=2n-\mult(1)-\mult(-1)$, and the same estimate
applies. (d) By (a) and (b), $B^+$ has a simple largest element and contains
every other value at most twice. Five elements need $\rho\ge4$ and
$S_2\ge4^2+2\cdot3^2+2\cdot2^2=42$, with equality only for $\{4,3,3,2,2\}$, and
six need $S_2\ge5^2+4^2+2\cdot3^2+2\cdot2^2=67$; and $S_2\le n$ by (c).
\end{proof}

\section{Odd cycles}\label{sec:odd}

\begin{proof}[Proof of Proposition~3]
For $v\ge2$ let $\eta_v=\mult(v)-\mult(-v)$. For odd $L$ the eigenvalues in
$\{0,\pm1\}$ satisfy $\lambda^L=\lambda$, and $(-v)^L-(-v)=-(v^L-v)$, so by
Lemma~\ref{lem:trace}(b) and $\tr A=0$,
\begin{equation}\label{eq:eta}
 \sum_{v\ge2}\eta_v\,(v^L-v)=\tr A^L=0\ \ (L\text{ odd},\ L<g),\qquad
 \sum_{v\ge2}\eta_v\,(v^g-v)=2g\,c_g.
\end{equation}
Let $q$ be prime with $(q-1)\mid(g-1)$. For every integer $\lambda$,
$\lambda^g\equiv\lambda\pmod q$: this is clear if $q\mid\lambda$, and otherwise
$\lambda^{q-1}\equiv1$ by Fermat's little theorem, so
$\lambda^g=\lambda(\lambda^{q-1})^{(g-1)/(q-1)}\equiv\lambda$. Hence $q$ divides
$\sum_\lambda(\lambda^g-\lambda)=2g\,c_g$. The prime $q=3$ qualifies as $g$ is
odd.

Let $V=\{v\ge2:\eta_v\ne0\}$, $s=|V|$, and suppose $s\le(g-3)/2$. The
exponents $L=2j+1$, $1\le j\le s$, are odd and less than $g$, so
$\sum_{v\in V}\eta_vv(w_v^j-1)=0$ for $1\le j\le s$, with $w_v=v^2$. Subtracting
the row of the value $1$ from the other rows of the Vandermonde matrix of the
distinct numbers $1$ and $w_v$ ($v\in V$) shows that $\det(w_v^j-1)$ equals
that Vandermonde determinant, which is nonzero. Hence all $\eta_v=0$, which
contradicts the definition of $V$ if $s>0$ and~\eqref{eq:eta} if $s=0$. So
$s\ge(g-1)/2$, and $2e=\sum_\lambda\lambda^2\ge\sum_{v\in V}v^2\ge
\sum_{v=2}^{(g+1)/2}v^2$.
\end{proof}

\begin{corollary}\label{cor:odd}
Let $G$ be an integral unicyclic graph with $n$ vertices whose cycle has odd
length $b$, and let $\kappa_b$ be the product of the primes $q$ with
$(q-1)\mid(b-1)$. Then $\kappa_b\mid2b$; in particular $3\mid b$. The odd
$b<200$ with $\kappa_b\mid2b$ are $3$, $15$, $27$, $39$, $63$, $75$, $87$,
$99$, $123$, $135$, $147$, $159$, $183$ and $195$. If $b\ne3$, then $b\ge15$ and
$n\ge102$.
\end{corollary}

\begin{proof}
Here $g=b$, $c_g=1$ and $e=n$, so $\kappa_b\mid2b$ by Proposition~3. For
$b=5,7,9,11,13$, $\kappa_b=30,42,30,66,2730$ does not divide $2b$; the list is a
direct computation. If $b\ge15$, then $2n\ge\sum_{v=2}^8v^2=203$.
\end{proof}

\begin{remark}\label{rem:novelty}
By the theorem of von Staudt and Clausen, $\kappa_b$ is the denominator of the
Bernoulli number $B_{b-1}$. The first part of Proposition~3 uses integrality;
for an arbitrary integer matrix one only has the classical congruence
$\tr A^q\equiv\tr A\pmod q$ for primes $q$. The second part is a variant of the
odd-polynomial argument of van Dam and Haemers~\cite[Section~2]{vdh}. We found
no earlier statement of Proposition~3, even for unicyclic graphs ($20$ web
searches, $9$ arXiv API queries and full-text searches of related papers), but
the paper of Omidi~\cite{om} on necessary conditions for integral unicyclic
graphs and~\cite{bdr} could not be read in full (paywalled). As the argument is
elementary, its novelty is uncertain; we present it as an observation.
\end{remark}

\section{Restricting the cycle length}\label{sec:cycle}

\begin{lemma}\label{lem:cycle}
$N_{>1}(G)\ge N_{>1}(C_b)=2\lceil b/6\rceil-1$ and
$N_{<-1}(G)\ge N_{<-1}(C_b)$. So if $G$ is integral and not a cycle, then
$|B^+|\ge2\lceil b/6\rceil-1$; if moreover $b$ is even, then $b\le12$ when
$n\le41$, and $b\le18$ when $n\le42$, with $B^+=\{4,3,3,2,2\}$ for
$b\in\{14,16,18\}$. In particular, if $n\le42$ and $b\equiv2\pmod 4$, then
$b\in\{6,10,14,18\}$.
\end{lemma}

\begin{proof}
$C_b$ is an induced subgraph of $G$, so by Cauchy interlacing~\cite{bh} the
$i$-th largest (smallest) eigenvalue of $G$ is at least (at most) that of $C_b$
for $i\le b$. The eigenvalues of $C_b$ are $2\cos(2\pi j/b)$, $j\in\Z_b$, and
$2\cos(2\pi j/b)>1$ if and only if $6|j|<b$ for $j\in(-b/2,b/2]$, which holds
for $2\lceil b/6\rceil-1$ values of $j$. If $G$ is integral, its eigenvalues
greater than $1$ form $B^+$. Now apply Lemma~\ref{lem:adm}(d);
$2\lceil b/6\rceil-1=5$ for $13\le b\le18$.
\end{proof}

\begin{lemma}\label{lem:count}
Let $G$ be integral with $b$ even, and $s=|B^+|$. Then
$\mult(\pm1)=n-S_2$ and $\mult(0)=2S_2-n-2s$. Hence $S_2\le n\le2S_2-2s$ and
$2S_2-n-2s\le\max(P,2)$.
\end{lemma}

\begin{proof}
By symmetry $n=\mult(0)+2\mult(1)+2s$, and~\eqref{eq:m4} gives
$2n=2S_2+2\mult(1)$. The last inequality is Lemma~\ref{lem:mult}(d).
\end{proof}

\begin{corollary}\label{cor:1418}
No integral generalized sun graph with cycle length $14$ or $18$ has at most
$42$ vertices.
\end{corollary}

\begin{proof}
By Lemma~\ref{lem:cycle}, $s=5$ and $S_2=42$. By Lemma~\ref{lem:count},
$42\le n\le42$, so $\mult(1)=0$ and $\mult(0)=84-42-10=32\le\max(P,2)$. Then
$n\ge b+P\ge46$, a contradiction.
\end{proof}

The next bound is due to the first referee agent (see the disclosure) and is the
pruning rule of implementation~C in Section~\ref{sec:comp}. For $t\ge1$ let
$\sigma(t)$ be the sum of the first $t$ terms of $4,4,4,4,9,9,9,9,16,\dots$, in
which each $v^2$, $v\ge2$, appears four times, and let $\sigma(t)=0$ for
$t\le0$. Thus $\sigma(9)=68$, $\sigma(10)=84$, $\sigma(11)=100$, $\sigma(14)=166$.

\begin{lemma}\label{lem:rank}
Let $G$ be integral and not a cycle, $r$ the number of eigenvalues with
$|\lambda|\ge2$, and $r_0=b+|K_p|-\delta_0$, with $\delta_0,\delta_{\pm1}$ as in
Lemma~\ref{lem:mult}. Then
\begin{itemize}
\item[(a)] $r=r_0+2|K_q|-\delta_1-\delta_{-1}\ge b-4$;
\item[(b)] $\sigma(r)\le2n$;
\item[(c)] $\sigma(r_0)\le2n-2|K_q|$ if $K_q\ne\emptyset$, and
 $\sigma(r_0-4)\le2n-4$ if $K_q=\emptyset$.
\end{itemize}
Consequently, $b\le13$ if $n\le42$.
\end{lemma}

\begin{proof}
(a) Since $r=n-\mult(0)-\mult(1)-\mult(-1)$, the formula follows from
Lemma~\ref{lem:mult}(b),(c). If $K_p\ne\emptyset$, then $\delta_0\le|K_p|$ and
$r_0\ge b$; otherwise $K_q\ne\emptyset$ and $r_0\ge b-2$. With
$\delta_{\pm1}\le|K_q|$ if $K_q\ne\emptyset$ and $\delta_{\pm1}\le2$ otherwise,
$r\ge b-4$. (b) Each value $\pm v$, $v\ge2$, occurs at most twice
(Lemma~\ref{lem:adm}(b)), so the squares of these $r$ eigenvalues sum to at
least $\sigma(r)$ and at most $2n$. (c) If $K_q\ne\emptyset$, put
$x=2|K_q|-\delta_1-\delta_{-1}\ge0$; then $r=r_0+x$ and
$\mult(1)+\mult(-1)=2Q-x$, so
$\sigma(r_0)+4x\le\sigma(r)\le2n-2Q+x$, whence $\sigma(r_0)\le2n-2|K_q|$. If
$K_q=\emptyset$, put $y=\delta_1+\delta_{-1}\le4$; then $r=r_0-y$ and
$\sigma(r_0-4)+4(4-y)\le\sigma(r_0-y)\le2n-y$ (trivially if $r_0<4$), whence
$\sigma(r_0-4)\le2n-4$.

Let $n\le42$. If $b\ge15$, (a) and (b) give $100=\sigma(11)\le84$. If $b=14$,
then $r_0\ge14$ (as $4\nmid14$), and (c) gives $166\le82$ or $84\le80$.
\end{proof}

\section{Computations}\label{sec:comp}

\subsection{The programs}
Each program takes $b$ and $N$, runs through the sequences $((p_k,q_k))_k$, not
all zero, with $n\le N$, allowing $p_k,q_k>0$ at the same vertex, and decides
exactly which graphs are integral. Graphs reported as integral were confirmed
by factoring the characteristic polynomial with SymPy~\cite{sympy}.

\emph{Implementation A} (\code{sunsearch.c}) lists the admissible spectra
(Lemma~\ref{lem:adm} and, for odd $b$, \eqref{eq:odd}; version~5 adds
Lemma~\ref{lem:cycle}, version~6 optionally Lemma~\ref{lem:count}) and assigns
$(p_k,q_k)$ to $v_0,v_1,\dots$ by depth-first search, with a lexicographically
largest pair at $v_0$. A node is discarded if no admissible spectrum is
compatible with the bounds of Lemma~\ref{lem:prune}, \eqref{eq:m4} and
Lemma~\ref{lem:adm}(c); as these bounds are monotone, the loops over $p_k$ and
$q_k$ stop at the first discarded value. With Lemma~\ref{lem:count} switched on,
version~6 also uses its inequalities in this test, with $n\ge b+P+2Q$ for the
partial sums $P$ and $Q$; they are monotone as well. The window of Lemma~\ref{lem:window}
only skips a node, and the last pair is solved from~\eqref{eq:m4}. At a leaf,
A computes $\mult(m)$ for all $|m|\le9$ by Lemma~\ref{lem:mult}, for $|m|\ge2$
with the integer matrices $hT_k$, $h=m(m^2-1)$, modulo four primes near
$2^{62}$ and with an a priori bound below $2^{240}$ that makes the tests exact;
$G$ is integral if and only if these multiplicities add up to $n$, as
$|\lambda|\le\sqrt{2n}<10$.

\emph{Implementation B} (\code{sunB.c}) was written separately, by the same
agent. Version~1 prunes
only by the largest admissible value of $\Theta(G)$, tests~\eqref{eq:m4} at the
leaves (using $\mult(\pm1)\le\max(Q,2)$), and then tests the dihedrally
canonical sequences (least image) exactly, via
\[
 F(x)=\det\bigl(x(x^2-1)(xI-A(C_b))-\diag(p_k(x^2-1)+q_kx^2)\bigr)
 =x^{b-P}(x^2-1)^{b-Q}\det(xI-A),
\]
the equality following from the proof of Lemma~\ref{lem:schur}. $F$ is monic of
degree $4b$, and if $G$ is integral its roots are integers in $[-9,9]$. B
computes $F$ modulo three primes by interpolation; if the multiplicities of
the roots $m\in[-9,9]$ modulo some prime add up to less than $4b$, then $G$ is
not integral. Version~2 adds Lemma~\ref{lem:cycle} and the window of
Lemma~\ref{lem:window}; version~3 adds Lemma~\ref{lem:count}, the rotation rule
of A and the greatest image as canonical form. \emph{B2} (\code{sunB2.c}) is B
(version~1) with the rotation rule of A, the greatest image as canonical form
and a floating-point interlacing test (LAPACK's \texttt{dsyev}, margin
$10^{-6}$) on an induced subgraph of every completion; version~2 adds
Lemma~\ref{lem:cycle} and takes a larger induced subgraph. B2 serves only as a
cross-check. (The header comments of \code{sunB.c}, version~3, and of both
versions of B2 still describe the least image.)

\emph{Implementation C} (\code{sunC.c}, \code{sunC2.c}) was written from scratch
by the first referee agent. It enumerates support patterns (whether $p_k>0$, whether
$q_k>0$) up to dihedral symmetry and then all values on each pattern. It
discards patterns by Lemma~\ref{lem:rank}(a),(b) (C), also by
Lemma~\ref{lem:rank}(c) (C2), or not at all. Its exact test computes
$\phi(G)=\det(xI-A)$ modulo two primes by Schwenk's formula~\cite{sch} for the
edge $e=v_{b-1}v_0$, $\phi(G)=\phi(G-e)-\phi(G-v_{b-1}-v_0)-2\phi(G-V(C_b))$,
and, as B does, adds up the root multiplicities of the integers in $[-10,10]$
modulo each prime, which must give $n$ if $G$ is integral.

\emph{Implementation D} (\code{sunD_v1.c}, \code{sunD2.c}, \code{sunD3.c}) was
written from scratch by a second referee agent (see the disclosure). It runs through all sequences,
with no rotation rule, or only through those that are canonical for the
greatest image under the key $64(p_k+2q_k)+p_k$. Its only pruning, for even
$b$, is exact reachability, computed by dynamic programming: some completion
and some multiset $B^+$ with a simple largest element $\rho\ge3$ and other
elements in $[2,\rho-1]$, each at most twice, must satisfy the second identity
of~\eqref{eq:m4} and $S_2\le n\le\min(N,S_2+\max(Q,2))$ (version~1 omits the
bound $S_2+\max(Q,2)$). Version~3 can switch the pruning off, and then also
accepts odd $b$. For each integer $k\in[-9,9]$, D bounds $\mult(k)$
from above by the nullity modulo $2^{31}-1$ of a $b\times b$ integer matrix:
$k(k^2-1)S(k)$ for $|k|\ge2$ (Lemma~\ref{lem:schur}), and the systems of
Lemma~\ref{lem:mult}(b),(c) for $k\in\{0,\pm1\}$. If these bounds add up to
less than $n$, then $G$ is not integral. The runs of D used no inertia bound and none of
Lemmas~\ref{lem:cycle}--\ref{lem:rank} or~\ref{lem:window}.

\subsection{Soundness of the pruning}
Let $X=\{0,\dots,b-2\}$. Below, a tridiagonal matrix indexed by $Y\subseteq X$
has off-diagonal entry $1$ in the positions $(k,k+1)$ and $(k+1,k)$ for which
both $k$ and $k+1$ lie in $Y$, and $0$ elsewhere.

\begin{lemma}\label{lem:prune}
Let $m\ge2$ and let $T$ be tridiagonal on $X$ with diagonal entries
$t_k\le f_k(m)-m$. Then $N_{>m}(G),N_{<-m}(G)\ge n_+(T)$ and
$N_{\ge m}(G),N_{\le-m}(G)\ge n_+(T)+n_0(T)$. If $Y\subseteq K_q\cap X$ and $T'$
is tridiagonal on $X\setminus Y$ with diagonal entries $t'_k\le p_k-1$, then
$N_{>1}(G),N_{<-1}(G)\ge|Y|+n_+(T')$.
\end{lemma}

\begin{proof}
The principal submatrix $S(m)_X$ is tridiagonal, as $b-1\notin X$, and
$S(m)_X-T$ is diagonal and nonnegative. By Weyl's inequalities and Cauchy
interlacing~\cite{bh}, $n_+(S(m))\ge n_+(S(m)_X)\ge n_+(T)$, and likewise for
$n_++n_0$; apply Lemma~\ref{lem:schur}(a). As $f_k$ is odd, $\Delta
S(-m)_X\Delta=-S(m)_X$ for $\Delta=\diag((-1)^k)$, and Lemma~\ref{lem:schur}(b)
gives the negative side. For $T'$, let $U=(K_q\cap X)\setminus Y$. By
Lemma~\ref{lem:schur}(c) and interlacing, $N_{>1}(G)\ge|K_q|+n_+(R_{Z\cap X})$;
and $R_{Z\cap X}$ arises from the tridiagonal matrix on $X\setminus Y$ with
diagonal $p_k-1$ by deleting the indices in $U$, so
$n_+(R_{Z\cap X})\ge n_+(T')-|U|$, while $|K_q|\ge|Y|+|U|$. The negative side
is the same with $R'$ and $\Delta$.
\end{proof}

At a node of A with $v_0,\dots,v_{d-1}$ assigned ($d\le b-1$), $T$ and $T'$ use
the assigned values and, elsewhere, those of a bare vertex ($t_k=-m$,
$t'_k=-1$), with $Y=K_q\cap\{0,\dots,d-1\}$. Raising $p_k$ or $q_k$ raises
$t_k,t'_k$ or moves $k$ into $Y$, which does not decrease $|Y|+n_+(T')$; so the
bounds are monotone. Inertias are computed by the pivot recurrence $d_1=t_1$,
$d_k=t_k-1/d_{k-1}$, restarted after each index missing from the index set, in
exact rational arithmetic (no overflow occurred).

\begin{lemma}\label{lem:pivot}
Let $T$ be symmetric tridiagonal with nonzero off-diagonal entries. If an exact
zero pivot $d_k=0$ is counted as negative, $d_{k+1}$ as positive, and the
recurrence is restarted with $d_{k+2}=t_{k+2}$, then the number of positive
pivots is $n_+(T)$; with the signs exchanged one obtains $n_-(T)$.
\end{lemma}

\begin{proof}
For small $\varepsilon>0$, $n_+(T)=n_+(T-\varepsilon I)$, whose pivots are
$\Delta_k(\varepsilon)/\Delta_{k-1}(\varepsilon)$ with leading minors $\Delta_k$.
The leading blocks are Jacobi matrices, with simple and strictly interlacing
eigenvalues; if $\Delta_k(0)=0$, counting their negative eigenvalues shows
that $\Delta_k(\varepsilon)$ and $\Delta_{k-1}(0)$ have opposite signs. The
recurrence $\Delta_{j+1}=t_{j+1}\Delta_j-a_j^2\Delta_{j-1}$ ($a_j\ne0$ the
off-diagonal entries) gives $\Delta_{k+1}(0)=-a_k^2\Delta_{k-1}(0)$ and
$\Delta_{k+2}(0)=t_{k+2}\Delta_{k+1}(0)$. So $d_k(\varepsilon)\to0^-$,
$d_{k+1}(\varepsilon)\to+\infty$, $d_{k+2}(\varepsilon)\to t_{k+2}$. For $n_-$
use $T+\varepsilon I$.
\end{proof}

\begin{lemma}\label{lem:window}
If $v_0,\dots,v_{d-1}$ are assigned, contributing $\Theta_d$ to $\Theta(G)$ and
$W$ to $n-b$, and $r=b-d$, then every completion with $n\le N$ has
$\Theta_d+2r\le\Theta(G)\le\Theta_d+g(N-b-W)+2(r-1)$, where $g(w)=(2+w)(1+w)$.
\end{lemma}

\begin{proof}
A vertex of weight $w=p+2q$ contributes $g(w-q)+2q\le g(w)$, as
$g(t)-g(t-1)=2t+2$, and at least $g(0)=2$. As $g$ is convex,
$g(a)+g(c)\le g(a+c)+g(0)$, and $g$ is increasing.
\end{proof}

A node is skipped if no admissible value of $\Theta(G)$ from~\eqref{eq:m4} lies
in this window. The window is not monotone in $p_k$, so it never ends a loop.

\subsection{Results, controls and versions}
Table~\ref{tab:runs} lists the searches; runs split into parts are summed.
Besides these, A (version~4, $15{,}678$ nodes), C (no pruning, $1{,}601{,}329$
sequences) and D (no pruning, all $3{,}031{,}469$ sequences) found no integral
generalized sun graph with cycle length $3$ and at most $49$ vertices.

\begin{table}[ht]\centering\footnotesize
\begin{tabular}{@{}cc>{\raggedright\arraybackslash}p{0.28\textwidth}>{\raggedright\arraybackslash}p{0.3\textwidth}>{\raggedright\arraybackslash}p{0.19\textwidth}@{}}\toprule
$N$ & $b$ & A & B, B2 & C; D\\\midrule
41 & 3 & v4: 5{,}033 / 1{,}357 & v1: 10{,}269 / 1{,}754 & 521{,}056 (np)\\
 & 4 & v4: 398{,}136 / 85{,}505 & v1: 410{,}544 / 52{,}641 & 5{,}683{,}267 (np);\newline D: 410{,}544 (all)\\
 & 5, 7, 9 & v4: no spectrum & v1: no spectrum & 20{,}292{,}905; 82{,}161{,}850;\newline 4{,}716{,}966\\
 & 8 & v4: 70{,}518{,}560 / 576{,}976 & & 57{,}561{,}471;\newline D: 2{,}156{,}331\\
 & 10 & v4: 50{,}716{,}095 / 52{,}652 & B2 v1: 52{,}652 / 14{,}662 & 3{,}814{,}116;\newline D: 2{,}124{,}182\\
 & 11, 13, 15 & & v1: no spectrum & $b=11$: 1{,}239{,}693; $b=13$: 0\\
 & 12 & v4: 5{,}860{,}901 / 902 & & 677{,}237;\newline D: 618{,}755\\\midrule
42 & 3 & v4: 5{,}334 / 1{,}381 & v1: 11{,}895 / 2{,}030;\newline v3: 3{,}991 / 2{,}030 & 606{,}367 (np);\newline D: 1{,}246{,}783 (np)\\
 & 6 & v4: 35{,}443{,}626 / 1{,}128{,}560;\newline v6: 11{,}623{,}493 / 572{,}032 & v1: 14{,}215{,}675 / 1{,}187{,}991;\newline v3: 1{,}046{,}629 / 520{,}217;\newline B2 v2: 1{,}128{,}560 / 449{,}500 & 201{,}673{,}369;\newline D: 153{,}225{,}129 (all)\\
 & 10 & v5: 465{,}028{,}393 / 461{,}049;\newline v6: 19{,}679{,}948 / 71{,}680 & v3: 401{,}175 / 200{,}579 & C2: 4{,}601{,}920;\newline D: 119{,}245{,}995 (all)\\
 & 14 & v5: 904{,}115 / 850 & v3: 0 & 1{,}951{,}786;\newline D: 1{,}128{,}742\\
 & 18 & v5: 18 / 0 & v2: 2{,}048{,}976 / 56{,}952;\newline v3 and B2 v2: 0 & D: 67{,}817\\
 & 5--15 odd & v4 ($b\le9$): no spectrum & v1, v3: no spectrum & \\
 & 22, 26 & v6: no spectrum & & D: 0\\\bottomrule
\end{tabular}
\caption{The searches with $n\le N$. For A: search-tree nodes / exact leaf
tests; for B and B2: sequences passing the moment tests / canonical sequences
tested exactly; for C and D: sequences tested exactly, (np) meaning without
pruning and (all) without symmetry reduction. The runs of D marked (np) have no
symmetry reduction either, and its other runs test canonical sequences only.
``No spectrum'': no admissible spectrum exists; for D, ``0'': no sequence passes
its pruning. Every run found nothing, except at $b=4$
(A, B, C and D: the three graphs of Theorem~2) and at $b=6$, $N=42$ (every run:
$C_{6,1}(0,6,6,12,6,6)$ only). Version~6 of A used Lemma~\ref{lem:count}.}
\label{tab:runs}
\end{table}

\paragraph{Controls and cross-checks.}
A and B also agree at $b=4$, $N=36$. Version~2 of B2 ran to completion at
$b=6$, $N=42$ and found only $C_{6,1}(0,6,6,12,6,6)$; the $1{,}128{,}560$
sequences passing its moment tests are as many as the exact leaf tests of A
(version~4). A note appended to its log says that this run was stopped, but
the log also has the final line, which B2 prints only on completion, and a
rerun gave the same final line. A run of B2 at $b=10$, $N=42$ was stopped
unfinished and is not used. Version~6 of A at $b=4$, $N=49$ gives the same list
with and without Lemma~\ref{lem:count}, including $C_{4,2}(10,10,0,0)$ with $44$
vertices and $\mult(\pm1)=19$, and so does D. The first referee reran
version~4 of A and version~1 of B2 at $b=10$, $N=41$: the sequences passing
their moment tests form the same set of $52{,}652$, none of them integral under
the exact test of C. Without symmetry reduction, D found at $b=6$, $N=42$
exactly the six sequences of the dihedral orbit of $C_{6,1}(0,6,6,12,6,6)$, and
nothing at $b=10$, $N=42$. The numbers of sequences passing its pruning and of
canonical sequences equal those of B (version~1) at $b=4$, $N=41$
($410{,}544$ and $52{,}641$) and at $b=6$, $N=42$ ($14{,}215{,}675$ and
$1{,}187{,}991$). Against exact characteristic polynomials (SymPy),
Lemma~\ref{lem:mult}, Lemma~\ref{lem:trace}(a) and the bound of
Lemma~\ref{lem:rank}(a) hold on $2{,}968$ random graphs ($b\le14$, $n\le60$;
$1{,}551$ with $\delta_0>0$, $953$ with $\delta_1+\delta_{-1}>0$), and the
polynomials of C agree modulo $2^{61}-1$ on $400$ graphs; the multiplicities of
A agree with floating-point eigenvalues on $3{,}877$ graphs, and the bounds of
D equal the exact multiplicities (nullities over $\mathbb Q$) on $3{,}250$
random graphs ($b\le12$, $n\le60$). The second referee also rebuilt versions~5
and~6 of A and version~3 of B from their hashed sources and reproduced the
logged counters of five of their runs. Two errors in earlier versions of A were found
by these controls and by self-review, and corrected before the runs of
Table~\ref{tab:runs}: the first version used the path on all $b$ cycle
vertices, not a principal submatrix, and missed $C_{6,1}(0,6,6,12,6,6)$;
version~3 ended loops when the window of Lemma~\ref{lem:window} failed. The
superseded logs are kept.

\paragraph{Versions.}
The programs, logs, test scripts and both referee reports are in
version~1.6.0 of~\cite{repo}, directory \code{certificates/sun-graphs/}, with
A and B in \code{code/}, C in \code{review/code/} and D in
\code{review/paper7/code/}; Table~\ref{tab:versions} lists the versions. The
first referee agent audited versions~4 of A and~1 of B and B2, which produced
all their runs with $N\le41$ and the runs of A and B with $N=42$ for
$b=3,5,6,7,9$ (A) and $b\le15$ (B). The later versions,
Lemma~\ref{lem:count} and Corollary~\ref{cor:1418} were written after its
report. The second referee agent then checked the two proofs line by line, read
the sources of versions~5 and~6 of A, versions~2 and~3 of B and version~2 of
B2, reproduced counters of the later versions as described above, and
recomputed the cases of Table~\ref{tab:cases} marked D with implementation~D.

\begin{table}[ht]\centering\footnotesize
\begin{tabular}{@{}ll>{\raggedright\arraybackslash}p{0.46\textwidth}@{}}\toprule
file & SHA-256, first 16 digits & runs\\\midrule
\code{sunsearch_v4_f5a38c28.c} & \texttt{f5a38c28cf395886} & A v4: $N\le41$; $N=42$, $b=3,5,6,7,9$; $N=49$, $b=3$\\
\code{sunsearch_v5_bfcca593.c} & \texttt{bfcca593dc3ea0b2} & A v5: $N=42$, $b=10,14,18$\\
\code{sunsearch_v6_39a87195.c} & \texttt{39a871959d2b3d92} & A v6: $N=42$, $b=6,10,22,26$; $N=49$, $b=4$\\
\code{sunB_v1_4d92eb48.c} & \texttt{4d92eb484cefea14} & B v1: $N\le42$, $b\le15$\\
\code{sunB_v2_8311d32e.c} & \texttt{8311d32e74f0cfbd} & B v2: $N=42$, $b=18$\\
\code{sunB.c} (version 3) & \texttt{f0eb3d9ca506a3f7} & B v3: $N=42$\\
\code{sunB2_v1_54531ad0.c} & \texttt{54531ad075d1015f} & B2 v1: $N=41$, $b=10$\\
\code{sunB2.c} (version 2) & \texttt{c80538f9d4359954} & B2 v2: $N=42$, $b=6,18$\\
\code{sunC.c} & \texttt{9401061cc17eff93} & C\\
\code{sunC2.c} & \texttt{b3e35c7bc97529ac} & C2: $N=42$, $b=10$\\
\code{sunD_v1.c} & \texttt{167d5a43514c4ef9} & D v1: $N=41$, $b=4$; $N=42$, $b=6$\\
\code{sunD2.c} & \texttt{5d913b626f9dc685} & D v2: $N=41$, $b=4,8$; $N=42$, $b=6,10$; $N=49$, $b=4$\\
\code{sunD3.c} & \texttt{e1a42832a46bd85d} & D v3: $N=41$, $b=10,12$; $N=42$, $b=14,18,22,26,30$\\
\code{sunD3_maxb16.c} & \texttt{18a172d9cb9e42da} & D v3, smaller arrays: $N=42,49$, $b=3$\\\bottomrule
\end{tabular}
\caption{Program versions. The full SHA-256 values are recorded with the logs.}
\label{tab:versions}
\end{table}

\section{Proofs of Theorems 1 and 2}\label{sec:proofs}

\begin{proof}[Proof of Theorem~1]
Let $G$ be a counterexample with $n\le42$; by~\cite[Theorem~2.1]{bms} it carries
only copies of $P_1$ and $P_2$. Odd $b\ge5$ is excluded by
Corollary~\ref{cor:odd}. If $b\equiv2\pmod4$, then $b\in\{6,10,14,18\}$ by
Lemma~\ref{lem:cycle}, and $b\notin\{14,18\}$ by Corollary~\ref{cor:1418}. So
$b\in\{3,6,10\}$, and the searches with $N=42$ give
$G\cong C_{6,1}(0,6,6,12,6,6)$, a counterexample by~\cite{bms}.
\end{proof}

\begin{proof}[Proof of Theorem~2]
Let $G$ be integral, not a cycle, with $n\le41$; by~\cite[Theorem~2.1]{bms} it
carries only copies of $P_1$ and $P_2$. By Corollary~\ref{cor:odd} and
Lemma~\ref{lem:cycle}, $b\in\{3,4,6,8,10,12\}$. The searches find nothing for
$b\in\{3,8,10,12\}$, only the graph on $42$ vertices for $b=6$, and exactly the
three graphs for $b=4$.
\end{proof}

\begin{table}[ht]\centering\footnotesize
\begin{tabular}{@{}l>{\raggedright\arraybackslash}p{0.36\textwidth}>{\raggedright\arraybackslash}p{0.4\textwidth}@{}}\toprule
case & by hand & by computation\\\midrule
\multicolumn{3}{@{}l}{\emph{Theorem~1: $n\le42$, $b\not\equiv0\pmod4$}}\\
$b=3$ & & A v4; B v1, v3; C and D (no pruning)\\
$b=6$ & & A v4, v6; B v1, v3; B2 v2; C; D\\
$b=10$ & & A v5 (without Lemma~\ref{lem:count}), v6; B v3; C2; D\\
$b=14$ & Corollary~\ref{cor:1418}; Lemma~\ref{lem:rank}(c) & A v5; C; D\\
$b=18$ & Corollary~\ref{cor:1418}; Lemma~\ref{lem:rank}(a),(b) & A v5; B v2; D\\
odd $b\ge5$ & Corollary~\ref{cor:odd}; for $b\ge15$ also Lemma~\ref{lem:rank}(a),(b) & A v4 ($b\le9$), B ($b\le15$): no spectrum\\
even $b\ge22$ & Lemma~\ref{lem:cycle}; Lemma~\ref{lem:rank}(a),(b) & A v6 ($b=22,26$): no spectrum; D ($b=22,26,30$): none\\\midrule
\multicolumn{3}{@{}l}{\emph{Theorem~2: $n\le41$}}\\
$b=4$ & & A v4; B v1; C (no pruning); D\\
$b=8,12$ & & A v4; C; D\\
$b=10$ & & also A v4; B2 v1; C and D at $N=41$\\
even $b\ge14$ & Lemma~\ref{lem:cycle}; Lemma~\ref{lem:rank}(a),(b) & \\
odd $b\ge5$ & Corollary~\ref{cor:odd}; for $b\ge15$ also Lemma~\ref{lem:rank}(a),(b) & C for $b\le13$, without Corollary~\ref{cor:odd}\\\bottomrule
\end{tabular}
\caption{Confirmations of each case. A and B, with B2, were written separately
by the agent that found the results; C, with C2, and D were written from
scratch by two referee agents. The cases $b=3,6,10$ of Theorem~2 are covered
by those of Theorem~1.}
\label{tab:cases}
\end{table}

\section{Open questions}\label{sec:open}
The questions of~\cite[Section~5]{bms}, whether the cycle length $b$ of an
integral generalized sun graph other than a cycle must be even, and whether $4$
and $6$ are the only possible values, remain open. By Corollary~\ref{cor:odd},
an odd $b$ is $3$, or at least $15$ with $n\ge102$; and $b=3$ needs $n\ge50$
(by the searches of Section~\ref{sec:comp}).
Lemma~\ref{lem:rank} bounds $b$ for either parity: if $r_{\max}(n)$ is the
largest $t$ with $\sigma(t)\le2n$, then $b\le r_{\max}(n)+4$; writing
$r_{\max}(n)=4k+j$ with $0\le j\le3$, we have
$2n\ge\sigma(4k)\ge4\sum_{v=1}^kv^2\ge\frac43k^3$, so $b\le(96n)^{1/3}+7$.

Thanks are due to Rodrigo O.\ Braga, Jean Carlo Moraes and Matheus C.\ Santos
for the question answered here.

\section*{Tool and computational resource disclosure}
This work was carried out on 26 September 2026 (UTC) in a workflow directed by
the author, using AI coding agents: Anthropic Claude Code, with the model
Claude Opus~5.5. Claude Code agents selected the problem, found the results and
their proofs, wrote the programs, carried out the literature searches
described in the introduction, and drafted this text. An independent Claude
Code agent, which had not produced the results, refereed the proofs line by
line, checked the lemmas of Section~\ref{sec:lemmas} against exact
computations, compared the leaf sets of two programs, and reran the searches
with its own program (implementation~C), written from scratch; its corrections
have been applied. A second referee agent, which had produced neither the
results nor implementation~C, later checked Lemma~\ref{lem:count},
Corollary~\ref{cor:1418} and the later program versions, and recomputed the
cases $b=3,4,6,8,10,12,14,18$ with a program of its own (implementation~D),
written from scratch (Section~\ref{sec:comp}); its corrections have been
applied as well. Throughout, ``we'' refers to this workflow; the
reading and searching described in the paper were done by the agents. The
searches ran on one computer (macOS, arm64), one core per run; no completed run
took more than an hour. The author is not a professional mathematician.
The author has read the paper.
Nothing in this paper is formally verified: the results rest on the written
proofs and on the computations of Section~\ref{sec:comp}. AI systems are not
authors of this paper; the author takes full responsibility for its content.

{\footnotesize
\begin{thebibliography}{99}\setlength{\itemsep}{1pt}
\bibitem{bdr}
R.~O. Braga, R.~R. Del-Vecchio and V.~M. Rodrigues, Integral unicyclic graphs,
\emph{Linear Algebra Appl.} \textbf{614} (2021), 281--300,
\url{https://doi.org/10.1016/j.laa.2020.05.002}.
\bibitem{bms}
R.~O. Braga, J.~C. Moraes and M.~C. Santos, Pendant paths and integral
generalized sun graphs, arXiv:2609.28754v1, 23 September 2026.
\bibitem{brt}
R.~O. Braga, V.~M. Rodrigues and V.~Trevisan, Locating eigenvalues of
unicyclic graphs, \emph{Appl.\ Anal.\ Discrete Math.} \textbf{11} (2017),
273--298, \url{https://doi.org/10.2298/AADM1702273B}.
\bibitem{bh}
A.~E. Brouwer and W.~H. Haemers, \emph{Spectra of Graphs}, Universitext,
Springer, New York, 2012, \url{https://doi.org/10.1007/978-1-4614-1939-6}.
\bibitem{repo}
H.~Dong, \emph{ai4math-results}, repository,
\url{https://github.com/Horace-Maxwell/ai4math-results}, version~1.6.0
(2026); archived at Zenodo under the concept DOI
\url{https://doi.org/10.5281/zenodo.22970254}.
\bibitem{hs}
F.~Harary and A.~J. Schwenk, Which graphs have integral spectra?, in
\emph{Graphs and Combinatorics}, Lecture Notes in Math.~406, Springer, Berlin,
1974, 45--51, \url{https://doi.org/10.1007/BFb0066434}.
\bibitem{hay}
E.~V. Haynsworth, Determination of the inertia of a partitioned Hermitian
matrix, \emph{Linear Algebra Appl.} \textbf{1} (1968), 73--81,
\url{https://doi.org/10.1016/0024-3795(68)90050-5}.
\bibitem{sympy}
A.~Meurer et al., SymPy: symbolic computing in Python, \emph{PeerJ Comput.\
Sci.} \textbf{3} (2017), e103, \url{https://doi.org/10.7717/peerj-cs.103}.
\bibitem{om}
G.~R. Omidi, On integral graphs with few cycles, \emph{Graphs Combin.}
\textbf{25} (2009), 841--849,
\url{https://doi.org/10.1007/s00373-010-0913-1}.
\bibitem{sch}
A.~J. Schwenk, Computing the characteristic polynomial of a graph, in
\emph{Graphs and Combinatorics}, Lecture Notes in Math.~406, Springer, Berlin,
1974, 153--172, \url{https://doi.org/10.1007/BFb0066438}.
\bibitem{tru}
The Omega Institute, \texttt{trureturing} repository,
\url{https://github.com/the-omega-institute/trureturing}, accessed 26
September 2026.
\bibitem{vdh}
E.~R. van Dam and W.~H. Haemers, An odd characterization of the generalized
odd graphs, \emph{J. Combin.\ Theory Ser.~B} \textbf{101} (2011), 486--489,
\url{https://doi.org/10.1016/j.jctb.2011.03.001}; also arXiv:1202.2300.
\bibitem{ver}
M.~S.~T. Veras, Buscando grafos unic\'{\i}clicos integrais [Searching for
integral unicyclic graphs], in \emph{XXXIII Sal\~ao de Inicia\c{c}\~ao
Cient\'{\i}fica}, Porto Alegre, Brazil, 2021 (in Portuguese),
\url{https://lume.ufrgs.br/bitstream/handle/10183/243963/Resumo_72921.pdf},
accessed 26 September 2026.
\end{thebibliography}}
\end{document}
